% Preamble template for Tufte-style lecture notes.
% Do not compile this file directly; compile the main notes file instead.
\documentclass{tufte-handout}

\usepackage{ntheorem}
\usepackage{graphicx}
\usepackage{amsmath}
\usepackage{amssymb}
\usepackage{hyperref}
\usepackage{epigraph}
\usepackage{booktabs}
\theoremstyle{break}
% \usepackage[
% bibencoding=utf8,% .bib file encoding
% maxbibnames=3, % otherwise et al
% minbibnames=1, % otherwise et al
% backend=biber,%
% sortlocale=en_US,%
% style=apa,% or authoryear
% % apabackref=false, % backreferences
% natbib=true,% for citet/citep, but this is for backward compatibility
% uniquename=false,%
% url=true,%
% sortcites=false,
% doi=true,%
% eprint=true%
% ]{biblatex}
% \addbibresource{lecture_note_bib.bib}

\hypersetup{
  colorlinks,
  urlcolor = blue,
  pdfauthor={Paul Goldsmith-Pinkham}
  pdfkeywords={econometrics}
  pdftitle={Lecture Notes for Applied Empirical Methods}
  pdfpagemode=UseNone
}
\newtheorem{ruleN}{Rule}
\newtheorem{thmN}{Theorem}
\newtheorem{assN}{Assumption}
\newtheorem{defN}{Definition}
\newtheorem{exmp}{Example}
\newtheorem{cmt}{Comment}
\newtheorem{proof}{Proof}

\newcommand{\continuation}{??}
\newtheorem*{excont}{Example \continuation}
\newenvironment{continueexample}[1]
 {\renewcommand{\continuation}{\ref{#1}}\excont[continued]}
 {\endexcont}
\newcommand{\bY}{\mathbf{Y}}
\newcommand{\bX}{\mathbf{X}}
\newcommand{\bD}{\mathbf{D}}
\newcommand{\E}{\mathbb{E}}

\newcommand\independent{\protect\mathpalette{\protect\independenT}{\perp}}
\def\independenT#1#2{\mathrel{\rlap{$#1#2$}\mkern2mu{#1#2}}}
\DeclareMathOperator{\Supp}{Supp}

\usepackage{cleveref}
\crefname{appsec}{appendix}{appendices}
\crefname{appsubsec}{appendix}{appendices}
\crefname{assumption}{assumption}{assumptions}
\crefname{equation}{equation}{equations}
\crefname{exmp}{example}{examples}
\crefname{assN}{assumption}{assumptions}
\crefname{cmt}{comment}{comments}
\crefname{defN}{definition}{definitions}

\usepackage[nolist]{acronym}
\begin{acronym}
  \acro{CI}{confidence interval}%
  \acro{OLS}{ordinary least squares}%
  \acro{CLT}{central limit theorem}%
  \acro{IV}{instrumental variables}%
  \acro{ATE}{average treatment effect}%
  \acro{RCT}{randomized control trial}%
  \acro{SUTVA}{stable unit treatment value assignment}
  \acro{VAM}{value-added model}%
  \acro{LAN}{locally asymptotically normal}%
  \acro{DiD}{difference-in-differences}%
  \acro{OVB}{omitted variables bias}
  \acro{FWL}{Frisch-Waugh-Lovell}
  \acro{DAG}{directed acyclic graph}
  \acro{PO}{potential outcomes}
\end{acronym}

\def\inprobHIGH{\,{\buildrel p \over \rightarrow}\,} 
\def\inprob{\,{\inprobHIGH}\,} 
\def\indistHIGH{\,{\buildrel d \over \rightarrow}\,} 
\def\indist{\,{\indistHIGH}\,}

\usepackage[many]{tcolorbox} 
\definecolor{main}{HTML}{3F6EA7}      % primary accent
\definecolor{sub}{HTML}{F2F7FD}       % cool background
\definecolor{sub2}{HTML}{FFF5E8}      % warm background
\definecolor{mainline}{HTML}{D3DFEE}  % soft frame for information boxes
\definecolor{warn}{HTML}{C7771B}      % accent for caution boxes
\definecolor{warnline}{HTML}{EAD9C0}  % soft frame for caution boxes

\tcbset{
    enhanced,
    breakable,
    colback = white,
    boxrule = 0.5pt,
    arc = 2pt,
    left = 9pt,
    right = 9pt,
    top = 7pt,
    bottom = 7pt,
    before skip = 0.3cm,
    after skip = 0.45cm
}                           % setting global options for tcolorbox

\newtcolorbox{boxD}{
    colback = sub,
    colframe = mainline,
    borderline west = {2.2pt}{0pt}{main},
    borderline north = {0.6pt}{0pt}{main!25},
    borderline south = {0.6pt}{0pt}{main!15}
}


\newtcolorbox{boxF}{
    colback = sub2,
    colframe = warnline,
    borderline west = {2.2pt}{0pt}{warn},
    borderline north = {0.6pt}{0pt}{warn!25},
    borderline south = {0.6pt}{0pt}{warn!15}
}

\usepackage{colortbl}


\usepackage{tikz}
\usepackage{verbatim}
\usetikzlibrary{positioning}
\usetikzlibrary{snakes}
\usetikzlibrary{calc}
\usetikzlibrary{arrows}
\usetikzlibrary{decorations.markings}
\usetikzlibrary{shapes.misc}
\usetikzlibrary{matrix,shapes,arrows,fit,tikzmark}




\title{S181M116 Derivatives and Alternative Investments: Lecture 5}
\author{Swapnil Singh}
\date{}

\hypersetup{
  pdftitle={S181M116 Derivatives and Alternative Investments: Lecture 5},
  pdfauthor={Swapnil Singh},
  pdfkeywords={derivatives, stochastic calculus, Brownian motion, Wiener process, Ito's lemma, geometric Brownian motion, lognormal distribution}
}

\setlength{\headheight}{26pt}

\newcommand{\dd}{\,d}
\newcommand{\dt}{\Delta t}
\newcommand{\dz}{\Delta z}
\newcommand{\eps}{\varepsilon}
\newcommand{\R}{\mathbb{R}}
\newcommand{\Prob}{\mathbb{P}}
\newcommand{\Var}{\operatorname{Var}}
\newcommand{\Cov}{\operatorname{Cov}}
\newcommand{\Corr}{\operatorname{Corr}}
\newcommand{\Normal}{\mathcal{N}}
\newcommand{\given}{\,|\,}
\newcommand{\St}{S_T}
\newcommand{\Szero}{S_0}
\newcommand{\Ito}{Ito}

\begin{document}
\maketitle

\begin{abstract}
This lecture develops the stochastic calculus used in continuous-time option
pricing. We begin with stochastic processes and the Markov property, then build
Brownian motion from independent normal increments. We add drift and volatility,
specialize the model to stock prices through geometric Brownian motion, and use
Monte Carlo simulation to interpret the model. The main mathematical tool is
\Ito's lemma, which explains how functions of stochastic variables evolve. The
final result is the lognormal distribution of stock prices, which is the input
for the Black-Scholes-Merton model in Lecture 6.
\end{abstract}

\begin{boxD}
\textbf{Reading.} Hull, Chapter 14. The goal is not to turn students into
probability theorists. The goal is to understand the stochastic notation well
enough that the Black-Scholes-Merton derivation in the next lecture feels like
economics plus no-arbitrage, not mysterious algebra.
\end{boxD}

\section{Roadmap}

Lecture 4 ended with binomial trees. A binomial tree is a discrete-time model:
over each short period the stock either moves up or down. Lecture 5 asks what
happens when the time step becomes very small and the stock is allowed to move
continuously.

The lecture has five blocks:
\begin{enumerate}
    \item \textbf{Stochastic processes and Markov dynamics.} Define the objects
    whose future is uncertain.
    \item \textbf{Brownian motion.} Build the basic source of continuous-time
    randomness.
    \item \textbf{Drift, volatility, and stock prices.} Explain why stock prices
    are modeled with percentage, not dollar, shocks.
    \item \textbf{Correlated processes and simulation.} Connect the math to
    practical Monte Carlo paths.
    \item \textbf{\Ito's lemma and lognormality.} Learn the chain rule of
    stochastic calculus and derive the distribution of future stock prices.
\end{enumerate}

\begin{boxF}
\textbf{Main lesson.} In ordinary calculus, small changes are of order $dt$.
In stochastic calculus, Brownian shocks are of order $\sqrt{dt}$. This is why
the square of a random shock matters: $(dz)^2$ behaves like $dt$ and creates
the extra $\frac{1}{2}$ term in \Ito's lemma.
\end{boxF}

\section{Stochastic Processes}

\subsection{What Is a Stochastic Process?}

A stochastic process is a collection of random variables indexed by time. If
$X_t$ is the process, then $X_t$ is the value of the variable at time $t$.

\begin{defN}[Stochastic process]
A stochastic process is a rule that assigns a random variable $X_t$ to each
date $t$ in a time set. The time set may be discrete, such as
$t=0,1,2,\ldots$, or continuous, such as $t\in[0,T]$.
\end{defN}

There are two separate distinctions.
\begin{description}
    \item[Discrete time versus continuous time] A process is observed or updated
    only at fixed dates in discrete time. In continuous time, changes can occur
    at any instant.
    \item[Discrete state versus continuous state] A process has a discrete state
    space if it can take only certain values. It has a continuous state space if
    it can take any value in an interval.
\end{description}

A binomial stock tree is discrete time and discrete state. A continuously
traded stock model is usually treated as continuous time and continuous state,
even though real exchanges trade in ticks and only during market hours.

\begin{boxD}
\textbf{Model versus reality.} Continuous-time models are idealizations. Stock
prices do not literally change at every mathematical instant, and quotes are
not infinitely divisible. The model is useful because it gives tractable
approximations and sharp no-arbitrage restrictions.
\end{boxD}

\subsection{The Markov Property}

\begin{defN}[Markov property]
A process has the Markov property if the conditional distribution of the future
depends on the current state, but not on the path taken to reach that state.
Formally, for $u>t$,
\[
    \Prob(X_u \leq x \given X_t, X_{t_1},\ldots,X_{t_n})
    =
    \Prob(X_u \leq x \given X_t),
\]
for all past dates $t_1,\ldots,t_n<t$.
\end{defN}

For a stock price, the Markov assumption says that if the stock is worth 100
now, then the distribution of its future price is summarized by the current
price and current model parameters. The fact that it was 90 last month or 120
last month does not directly matter after we condition on today's price.

\begin{boxF}
\textbf{Economic intuition.} A Markov stock price is the continuous-time cousin
of weak-form market efficiency. If past price patterns contained reliable
profit opportunities, traders would act on them, and their trades would tend to
move the information into today's price.
\end{boxF}

The Markov property does not say that past data are useless. Historical prices
can help estimate parameters such as volatility. It says that, after the model
parameters are known, the current state is enough for forecasting the future
distribution.

\section{Variance and the Square Root of Time}

\subsection{Adding Independent Normal Changes}

Suppose a Markov variable has a normally distributed one-year change with mean
zero and variance one:
\[
    \Delta X_{\text{1 year}} \sim \Normal(0,1).
\]
If annual changes are independent, the two-year change is the sum of two
independent normal variables. Means add and variances add:
\[
    \Delta X_{\text{2 years}} \sim \Normal(0,2).
\]
The standard deviation is therefore $\sqrt{2}$, not $2$.

More generally, if the variance rate is one per year, then over a time interval
of length $T$ years,
\[
    \Delta X_T \sim \Normal(0,T),
    \qquad
    \operatorname{sd}(\Delta X_T)=\sqrt{T}.
\]

\begin{boxD}
\textbf{Square-root-time rule.} Variance scales linearly with time. Standard
deviation scales with the square root of time. This is why daily volatility is
annual volatility divided by about $\sqrt{252}$, and annual volatility is daily
volatility multiplied by about $\sqrt{252}$.
\end{boxD}

\begin{exmp}[Scaling uncertainty]
If a process has annual standard deviation $20\%$, then under the square-root
rule its one-month standard deviation is approximately
\[
    20\%\sqrt{\frac{1}{12}} = 5.77\%.
\]
Its one-week standard deviation, using 52 weeks per year, is
\[
    20\%\sqrt{\frac{1}{52}} = 2.77\%.
\]
The one-month uncertainty is not one twelfth of the annual uncertainty because
standard deviations do not add across independent periods.
\end{exmp}

\section{Wiener Processes}

\subsection{Definition}

A Wiener process, also called standard Brownian motion, is the basic source of
randomness in the continuous-time stock price model.

\begin{defN}[Wiener process]
A process $z_t$ is a Wiener process if:
\begin{enumerate}
    \item for a small interval $\dt$, the increment is
    \[
        \dz = \eps \sqrt{\dt},
        \qquad
        \eps\sim\Normal(0,1);
    \]
    \item increments over non-overlapping time intervals are independent.
\end{enumerate}
\end{defN}

Thus,
\[
    \E[\dz]=0,
    \qquad
    \Var(\dz)=\dt,
    \qquad
    \operatorname{sd}(\dz)=\sqrt{\dt}.
\]
Over a longer interval $T$,
\[
    z_T-z_0 \sim \Normal(0,T).
\]

\subsection{Why Brownian Paths Are Jagged}

Brownian increments are proportional to $\sqrt{\dt}$. For small $\dt$,
$\sqrt{\dt}$ is much larger than $\dt$. For example, if
$\dt=0.0001$, then $\sqrt{\dt}=0.01$. The random term is therefore large
relative to deterministic time drift over very short intervals.

This is the source of the strange-looking Brownian path. It is continuous, but
it is nowhere differentiable in the ordinary sense. There is no well-defined
instantaneous slope.

\begin{boxF}
\textbf{Intuition.} A smooth curve has small movements of order $dt$. Brownian
motion has small movements of order $\sqrt{dt}$. As the time grid becomes
finer, the path keeps revealing new short-run variation instead of settling
down to a smooth tangent line.
\end{boxF}

\subsection{Differential Notation}

In stochastic calculus, we write
\[
    dz
\]
for the infinitesimal increment of a Wiener process. This notation should be
read through its discrete approximation:
\[
    dz \approx \eps \sqrt{dt}.
\]
The notation is compact, but the scale is important. The object $dz$ is random
and has variance $dt$.

\begin{exmp}[Distribution of a Brownian value]
Suppose $z_0=25$. Since $z_T-z_0\sim\Normal(0,T)$,
\[
    z_T \sim \Normal(25,T).
\]
At $T=1$, the standard deviation is $1$. At $T=5$, the standard deviation is
$\sqrt{5}=2.236$.
\end{exmp}

\section{Generalized Wiener Processes}

\subsection{Adding Drift and Scale}

The standard Wiener process has zero drift and variance rate one. A generalized
Wiener process adds a deterministic drift and rescales the noise:
\[
    dx = a\,dt + b\,dz,
\]
where $a$ and $b$ are constants.

Over a small time interval,
\[
    \Delta x = a\dt + b\eps\sqrt{\dt}.
\]
Therefore,
\[
    \E[\Delta x]=a\dt,
    \qquad
    \Var(\Delta x)=b^2\dt.
\]
Over a longer interval $T$,
\[
    x_T-x_0\sim\Normal(aT,b^2T).
\]

\begin{defN}[Drift rate and variance rate]
In $dx=a\,dt+b\,dz$, the drift rate is $a$ and the variance rate is $b^2$.
The diffusion coefficient $b$ determines the standard deviation per square root
of time.
\end{defN}

\begin{exmp}[Cash position]
A company's cash balance, measured in thousands, follows
\[
    dx = 20\,dt + 30\,dz.
\]
The drift rate is 20 thousand per year and the variance rate is
$30^2=900$. If $x_0=50$, then after one year
\[
    x_1 \sim \Normal(70,900),
\]
so the mean is 70 and the standard deviation is 30. After six months,
\[
    x_{0.5}\sim\Normal(60,450),
\]
so the standard deviation is $\sqrt{450}=21.21$.
\end{exmp}

\subsection{What Drift and Diffusion Do}

The drift term $a\,dt$ moves the expected path. The diffusion term $b\,dz$
creates dispersion around that path. In a very short interval, the stochastic
term is usually the visually dominant part because it is of order $\sqrt{dt}$.
Over longer horizons, the drift accumulates linearly while uncertainty grows
with the square root of time.

\begin{boxD}
\textbf{Useful mental picture.} Drift controls the direction of the cloud's
center. Volatility controls the width of the cloud. A single realized path can
move against the drift for a long time because the diffusion term is random.
\end{boxD}

\section{\Ito Processes}

\subsection{State-Dependent Drift and Volatility}

A generalized Wiener process has constant $a$ and $b$. An \Ito process allows
them to depend on the current state and time:
\[
    dx = a(x,t)\,dt + b(x,t)\,dz.
\]
Over a small interval,
\[
    \Delta x
    \approx
    a(x,t)\dt + b(x,t)\eps\sqrt{\dt}.
\]

This is a local approximation. During the short interval from $t$ to $t+\dt$,
we evaluate drift and volatility at the current state $(x,t)$.

\begin{boxF}
\textbf{Local thinking.} Continuous-time finance is local. We ask what happens
over the next instant given the current state, then aggregate the implications.
This is why derivatives can be priced by local hedging arguments.
\end{boxF}

If $a(x,t)$ and $b(x,t)$ depend only on current $x$ and $t$, the process is
Markov. If they depend on the entire past path, the process is generally not
Markov.

\section{The Stock Price Process}

\subsection{Why Not Constant Dollar Volatility?}

It is tempting to model a stock as
\[
    dS = a\,dt + b\,dz.
\]
This says the expected dollar change and dollar volatility are constant. That
is not natural for stocks. A 2 dollar daily standard deviation is huge for a
10 dollar stock and small for a 200 dollar stock.

Stock investors usually think in percentage returns. If required expected
return is 10 percent per year, that means 10 percent whether the stock price is
20 or 200. Likewise, volatility is usually quoted as a percentage of price.

\subsection{Geometric Brownian Motion}

The standard continuous-time stock price model is
\[
    dS = \mu S\,dt + \sigma S\,dz,
\]
or equivalently
\[
    \frac{dS}{S} = \mu\,dt + \sigma\,dz.
\]

\begin{defN}[Geometric Brownian motion]
A stock price follows geometric Brownian motion if
\[
    \frac{dS}{S} = \mu\,dt + \sigma\,dz,
\]
where $\mu$ is the expected continuously compounded rate of return and
$\sigma$ is the volatility.
\end{defN}

The proportional change over a short interval is approximately
\[
    \frac{\Delta S}{S}
    =
    \mu\dt+\sigma\eps\sqrt{\dt}.
\]
Thus,
\[
    \frac{\Delta S}{S}
    \sim
    \Normal(\mu\dt,\sigma^2\dt)
    \quad
    \text{approximately for small } \dt.
\]

\begin{boxD}
\textbf{Interpretation of $\mu$ and $\sigma$.} The parameter $\mu$ moves the
expected return. The parameter $\sigma$ is the annualized standard deviation of
continuously compounded returns. In derivative pricing, $\sigma$ is central.
The real-world $\mu$ will disappear from the no-arbitrage price in Lecture 6.
\end{boxD}

\begin{exmp}[Weekly stock return approximation]
A stock has expected return $\mu=15\%$ and volatility $\sigma=30\%$ per year.
For one week, take $\dt=1/52=0.01923$. Then
\[
    \frac{\Delta S}{S}
    =
    0.15(0.01923)+0.30\eps\sqrt{0.01923}
    =
    0.00288+0.0416\eps.
\]
The expected weekly return is about $0.288\%$ and the weekly standard deviation
is about $4.16\%$.
\end{exmp}

\subsection{Why the Model Is Called Geometric}

Without uncertainty,
\[
    dS=\mu S\,dt
\]
has solution
\[
    S_T=S_0e^{\mu T}.
\]
The stock grows geometrically, not linearly. With uncertainty, the same
multiplicative idea remains: shocks affect percentage changes, so prices tend
to stay positive and compound over time.

\section{Monte Carlo Simulation}

\subsection{Euler Simulation}

The simplest simulation uses the discrete approximation
\[
    S_{t+\dt}
    =
    S_t + \mu S_t\dt + \sigma S_t\eps\sqrt{\dt}
    =
    S_t\left(1+\mu\dt+\sigma\eps\sqrt{\dt}\right).
\]
At each step we draw a new independent $\eps\sim\Normal(0,1)$.

\begin{exmp}[One simulated step]
Let $S_t=100$, $\mu=0.15$, $\sigma=0.30$, $\dt=1/52$, and suppose the normal
draw is $\eps=0.52$. Then
\[
    \Delta S
    =
    100(0.00288+0.0416(0.52))
    =
    2.45.
\]
The simulated next price is $102.45$.
\end{exmp}

\subsection{Exact Lognormal Simulation}

For geometric Brownian motion with constant $\mu$ and $\sigma$, we can simulate
more accurately using the exact transition:
\[
    S_{t+\dt}
    =
    S_t
    \exp\left[
        \left(\mu-\frac{1}{2}\sigma^2\right)\dt
        +\sigma\sqrt{\dt}\eps
    \right].
\]
This formula comes from applying \Ito's lemma to $\ln S$, derived below.

\begin{boxF}
\textbf{Practical point.} Euler simulation is easy and useful for intuition.
Exact lognormal simulation is preferred for geometric Brownian motion because
it preserves positivity and avoids time-step approximation error for this
specific model.
\end{boxF}

\section{Volatility and Parameters}

\subsection{Annualization}

If $\sigma$ is annual volatility, then the standard deviation of the
continuously compounded return over $T$ years is
\[
    \sigma\sqrt{T}.
\]
For one trading day with 252 trading days in a year, it is
\[
    \frac{\sigma}{\sqrt{252}}.
\]
Conversely, if daily log-return standard deviation is $s_d$, annual volatility
is approximately
\[
    s_d\sqrt{252}.
\]

\begin{exmp}[Daily and annual volatility]
If annual volatility is $32\%$, daily volatility is approximately
\[
    \frac{0.32}{\sqrt{252}}=0.0202,
\]
or $2.02\%$. If daily volatility is $1.4\%$, annual volatility is approximately
\[
    0.014\sqrt{252}=0.222,
\]
or $22.2\%$.
\end{exmp}

\subsection{Why $\mu$ Is Harder Than $\sigma$}

Expected returns are hard to estimate because the signal is small relative to
noise. Volatility is easier because squared returns reveal the scale of
fluctuations even when average returns are imprecise.

For derivative pricing, this distinction matters. In the Black-Scholes-Merton
model, the derivative price depends on volatility, not on the stock's real-world
expected return. The reason is dynamic hedging: the stock's risk premium affects
both the derivative and the hedge in a way that cancels out.

\section{Correlated Wiener Processes}

\subsection{Two Sources of Risk}

Many derivatives depend on more than one risk factor. Examples include options
on exchange rates, spread options, quanto options, and derivatives with
stochastic volatility or stochastic interest rates.

Suppose
\[
    dx_1=a_1\,dt+b_1\,dz_1,
    \qquad
    dx_2=a_2\,dt+b_2\,dz_2.
\]
If the Brownian shocks have correlation $\rho$, then
\[
    \Corr(dz_1,dz_2)=\rho.
\]
In differential notation this is often summarized by the multiplication rule
\[
    dz_1\,dz_2 = \rho\,dt.
\]

\subsection{Simulating Correlated Normals}

To simulate two correlated standard normal shocks, draw independent
$u,v\sim\Normal(0,1)$ and set
\[
    \eps_1=u,
    \qquad
    \eps_2=\rho u+\sqrt{1-\rho^2}\,v.
\]
Then $\eps_1$ and $\eps_2$ each have standard normal distributions and their
correlation is $\rho$.

\begin{exmp}[Constructing correlated shocks]
Let $\rho=0.60$. Suppose $u=1.10$ and $v=-0.40$. Then
\[
    \eps_1=1.10,
    \qquad
    \eps_2=0.60(1.10)+\sqrt{1-0.60^2}(-0.40).
\]
Since $\sqrt{1-0.36}=0.80$,
\[
    \eps_2=0.66-0.32=0.34.
\]
The second shock shares a common component with the first shock but also has
independent residual variation.
\end{exmp}

\section{The Rules of Stochastic Calculus}

\subsection{Quadratic Variation}

The defining scale of Brownian motion is
\[
    dz \sim \sqrt{dt}.
\]
Therefore,
\[
    (dz)^2 \sim dt.
\]
This is the key difference from ordinary calculus. Terms involving $(dt)^2$ and
$dt\,dz$ are negligible, but $(dz)^2$ is not negligible because it is of order
$dt$.

The informal multiplication table is
\[
    (dt)^2=0,
    \qquad
    dt\,dz=0,
    \qquad
    (dz)^2=dt.
\]
For correlated Brownian motions,
\[
    dz_i\,dz_j=\rho_{ij}\,dt.
\]

\begin{boxF}
\textbf{Why the table works.} In a small interval, $dt$ is tiny and $dz$ is of
size $\sqrt{dt}$. Thus $(dt)^2$ is much smaller than $dt$, and $dt\,dz$ is of
size $dt^{3/2}$, also much smaller than $dt$. But $(dz)^2$ is of size $dt$ and
survives in the limit.
\end{boxF}

\subsection{Taylor Expansion With a Random Term}

Let $x$ follow
\[
    dx=a(x,t)\,dt+b(x,t)\,dz.
\]
Let $G(x,t)$ be a smooth function. The second-order Taylor expansion is
\[
    dG
    =
    G_x\,dx+G_t\,dt
    +\frac{1}{2}G_{xx}(dx)^2
    +G_{xt}\,dx\,dt
    +\frac{1}{2}G_{tt}(dt)^2
    +\cdots.
\]
Using the stochastic multiplication table, the only second-order term that
survives is $(dx)^2$:
\[
    (dx)^2
    =
    (a\,dt+b\,dz)^2
    =
    b^2(dz)^2
    =
    b^2dt.
\]
Therefore,
\[
    dG
    =
    G_x(a\,dt+b\,dz)+G_t\,dt+\frac{1}{2}G_{xx}b^2dt.
\]

\section{\Ito's Lemma}

\begin{thmN}[\Ito's lemma, one state variable]
If
\[
    dx=a(x,t)\,dt+b(x,t)\,dz,
\]
and $G=G(x,t)$ is sufficiently smooth, then
\[
    dG
    =
    \left(
        G_xa+G_t+\frac{1}{2}G_{xx}b^2
    \right)dt
    +
    G_xb\,dz.
\]
\end{thmN}

The new term relative to ordinary calculus is
\[
    \frac{1}{2}G_{xx}b^2\,dt.
\]
It appears because the stochastic variable has nonzero quadratic variation.

\begin{boxD}
\textbf{Financial interpretation.} If $G$ is an option price and $x$ is the
stock price, then $G_x$ is delta and $G_{xx}$ is gamma. \Ito's lemma says that
option value changes because of delta exposure, direct passage of time, and
curvature interacting with variance.
\end{boxD}

\subsection{Example: Squaring a Brownian Motion}

Let $x=z$ be a Wiener process, so $dx=dz$. Take
\[
    G(x)=x^2.
\]
Then $G_x=2x$, $G_{xx}=2$, and $G_t=0$. Since $a=0$ and $b=1$,
\Ito's lemma gives
\[
    d(x^2)=\frac{1}{2}(2)(1)\,dt+2x\,dz
    =
    dt+2x\,dz.
\]

Ordinary calculus would give only $2x\,dx$. Stochastic calculus adds $dt$.
That extra term is the accumulated variance of Brownian motion.

\subsection{Example: Exponential of Brownian Motion}

Let
\[
    dx=\mu\,dt+\sigma\,dz,
    \qquad
    G=e^x.
\]
Here $G_x=G$ and $G_{xx}=G$. \Ito's lemma gives
\[
    dG
    =
    G\left(\mu+\frac{1}{2}\sigma^2\right)dt
    +
    G\sigma\,dz.
\]
The exponential transformation changes the drift by
$\frac{1}{2}\sigma^2$. Convexity plus randomness raises the expected change in
the exponential.

\section{Applying \Ito's Lemma to Stock Prices}

\subsection{General Function of a Stock Price}

If
\[
    dS=\mu S\,dt+\sigma S\,dz,
\]
and $G=G(S,t)$, then $a=\mu S$ and $b=\sigma S$. \Ito's lemma becomes
\[
    dG
    =
    \left(
        G_S\mu S+G_t+\frac{1}{2}G_{SS}\sigma^2S^2
    \right)dt
    +
    G_S\sigma S\,dz.
\]

The same $dz$ appears in both $dS$ and $dG$. This common source of uncertainty
is what allows a derivative to be hedged locally using the underlying stock.

\begin{boxF}
\textbf{Bridge to Lecture 6.} If an option and the underlying stock are driven
by the same Brownian shock, a portfolio of the option and the stock can be
chosen so that the $dz$ terms cancel. The resulting locally riskless portfolio
must earn the risk-free rate. That is the engine behind the Black-Scholes-
Merton differential equation.
\end{boxF}

\subsection{Application: Forward Price}

For a non-dividend-paying stock with constant interest rate $r$, the forward
price for maturity $T$ at time $t$ is
\[
    F(S,t)=S e^{r(T-t)}.
\]
The derivatives are
\[
    F_S=e^{r(T-t)},
    \qquad
    F_{SS}=0,
    \qquad
    F_t=-rS e^{r(T-t)}.
\]
Using \Ito's lemma,
\[
    dF
    =
    \left(\mu S e^{r(T-t)}-rS e^{r(T-t)}\right)dt
    +
    \sigma S e^{r(T-t)}dz.
\]
Since $F=S e^{r(T-t)}$,
\[
    dF=(\mu-r)F\,dt+\sigma F\,dz.
\]

\begin{boxD}
\textbf{Interpretation.} The forward price has the same volatility as the stock
because it is proportional to the stock price. Its real-world drift is
$\mu-r$ because the forward price already compounds the spot price forward at
rate $r$.
\end{boxD}

\section{The Lognormal Property}

\subsection{Applying \Ito's Lemma to $\ln S$}

Let
\[
    G(S)=\ln S.
\]
Then
\[
    G_S=\frac{1}{S},
    \qquad
    G_{SS}=-\frac{1}{S^2},
    \qquad
    G_t=0.
\]
With $dS=\mu S\,dt+\sigma S\,dz$,
\[
    d\ln S
    =
    \left[
        \frac{1}{S}\mu S
        +
        \frac{1}{2}\left(-\frac{1}{S^2}\right)\sigma^2S^2
    \right]dt
    +
    \frac{1}{S}\sigma S\,dz.
\]
Therefore,
\[
    d\ln S
    =
    \left(\mu-\frac{1}{2}\sigma^2\right)dt
    +
    \sigma\,dz.
\]

\begin{boxF}
\textbf{The Ito correction.} The expected log return is not $\mu$. It is
$\mu-\frac{1}{2}\sigma^2$. Volatility lowers expected log growth because losses
hurt compounded wealth more than equal-sized gains help it.
\end{boxF}

\subsection{Distribution of Future Stock Prices}

Integrating from 0 to $T$,
\[
    \ln S_T-\ln S_0
    =
    \left(\mu-\frac{1}{2}\sigma^2\right)T+\sigma(z_T-z_0).
\]
Since $z_T-z_0\sim\Normal(0,T)$,
\[
    \ln S_T
    \sim
    \Normal\left(
        \ln S_0+\left(\mu-\frac{1}{2}\sigma^2\right)T,
        \sigma^2T
    \right).
\]
Equivalently,
\[
    S_T
    =
    S_0
    \exp\left[
        \left(\mu-\frac{1}{2}\sigma^2\right)T
        +
        \sigma\sqrt{T}\eps
    \right],
    \qquad
    \eps\sim\Normal(0,1).
\]

\begin{defN}[Lognormal distribution]
A positive random variable is lognormally distributed if its natural logarithm
is normally distributed. Under geometric Brownian motion, $S_T$ is lognormal.
\end{defN}

\subsection{Mean, Median, and the Volatility Drag}

The expected stock price under the real-world process is
\[
    \E[S_T]=S_0e^{\mu T}.
\]
The median stock price is
\[
    \operatorname{median}(S_T)
    =
    S_0e^{(\mu-\frac{1}{2}\sigma^2)T}.
\]
The mean exceeds the median because the lognormal distribution is right-skewed:
rare large upward outcomes pull up the average.

\begin{exmp}[Distribution after one year]
Let $S_0=100$, $\mu=12\%$, $\sigma=30\%$, and $T=1$. Then
\[
    \ln S_1
    \sim
    \Normal\left(
        \ln 100+(0.12-0.045),0.09
    \right).
\]
The mean log return is $7.5\%$ and the standard deviation of the log return is
$30\%$. The expected stock price is
\[
    100e^{0.12}=112.75,
\]
while the median stock price is
\[
    100e^{0.075}=107.79.
\]
\end{exmp}

\section{From Real-World to Risk-Neutral Dynamics}

Lecture 5 is mainly about stochastic calculus, but it is useful to preview the
pricing move. Under the real-world probability measure, the stock model is
\[
    \frac{dS}{S}=\mu\,dt+\sigma\,dz.
\]
Under the risk-neutral measure used for pricing a non-dividend-paying stock,
the drift becomes the risk-free rate:
\[
    \frac{dS}{S}=r\,dt+\sigma\,dz^{Q}.
\]
The volatility is unchanged in the basic Black-Scholes-Merton model.

\begin{boxD}
\textbf{Do not over-interpret risk-neutral probabilities.} Risk-neutral
probabilities are pricing probabilities, not forecasts. They are constructed so
that discounted asset prices are martingales and expected payoffs can be
discounted at the risk-free rate.
\end{boxD}

The corresponding risk-neutral terminal distribution is
\[
    \ln S_T
    \sim
    \Normal\left(
        \ln S_0+\left(r-\frac{1}{2}\sigma^2\right)T,
        \sigma^2T
    \right).
\]
This is the distribution used to derive the Black-Scholes-Merton formula.

\section{Common Mistakes}

\begin{enumerate}
    \item \textbf{Adding standard deviations across time.} Variances add;
    standard deviations scale with the square root of time.
    \item \textbf{Treating $dz$ like $dt$.} Brownian increments are of order
    $\sqrt{dt}$, so $(dz)^2$ contributes a $dt$ term.
    \item \textbf{Forgetting the \Ito correction.} If $dS/S=\mu dt+\sigma dz$,
    then $d\ln S=(\mu-\sigma^2/2)dt+\sigma dz$.
    \item \textbf{Confusing arithmetic and log returns.} The expected
    arithmetic instantaneous return is $\mu$, but expected log growth is
    $\mu-\sigma^2/2$.
    \item \textbf{Using real-world probabilities for pricing.} Option pricing
    uses risk-neutral dynamics after no-arbitrage arguments, not the stock's
    actual expected return.
\end{enumerate}

\section{Optional Extension: Fractional Brownian Motion}

Standard Brownian motion has independent increments and variance proportional
to elapsed time:
\[
    \Var(z_t-z_s)=t-s.
\]
Fractional Brownian motion generalizes this to
\[
    \Var(X_t-X_s)=\sigma^2(t-s)^{2H},
\]
where $H$ is the Hurst exponent. When $H=1/2$, the usual Brownian scaling is
recovered. When $H\neq 1/2$, increments are no longer independent in the same
way. Such models are useful in modern rough-volatility work, but the classical
Black-Scholes-Merton model uses ordinary Brownian motion.

\section{Worked Problems}

\begin{exmp}[From annual to quarterly uncertainty]
A stock has volatility $24\%$ per year. The standard deviation of its
continuously compounded return over one quarter is
\[
    0.24\sqrt{0.25}=0.12.
\]
Thus the quarterly log-return standard deviation is $12\%$.
\end{exmp}

\begin{exmp}[Probability of a log return event]
Let $S_0=50$, $\mu=8\%$, $\sigma=20\%$, and $T=1$. What is
$\Prob(S_T>55)$ under the real-world process?

Since
\[
    \ln S_T
    \sim
    \Normal\left(\ln 50+(0.08-0.5(0.20)^2),0.20^2\right),
\]
we have
\[
    \ln S_T
    \sim
    \Normal(\ln 50+0.06,0.04).
\]
Therefore
\[
    \Prob(S_T>55)
    =
    \Prob\left(
        Z>
        \frac{\ln(55)-\ln(50)-0.06}{0.20}
    \right).
\]
The numerator is approximately $0.0953-0.06=0.0353$, so the $z$-score is
$0.1765$. The probability is about $1-N(0.1765)=0.43$.
\end{exmp}

\begin{exmp}[\Ito's lemma for a power payoff]
Let $G(S)=S^n$ and
\[
    dS=\mu S\,dt+\sigma S\,dz.
\]
Then
\[
    G_S=nS^{n-1},
    \qquad
    G_{SS}=n(n-1)S^{n-2}.
\]
\Ito's lemma gives
\[
    dG
    =
    \left[
        n\mu S^n
        +
        \frac{1}{2}n(n-1)\sigma^2S^n
    \right]dt
    +
    n\sigma S^n dz.
\]
Since $G=S^n$,
\[
    \frac{dG}{G}
    =
    \left[
        n\mu+\frac{1}{2}n(n-1)\sigma^2
    \right]dt
    +
    n\sigma dz.
\]
Convex power payoffs have drift terms that increase with variance.
\end{exmp}

\begin{exmp}[Risk-neutral lognormal distribution]
A non-dividend-paying stock has $S_0=80$, $r=5\%$, $\sigma=25\%$, and
$T=0.5$. Under the risk-neutral measure,
\[
    \ln S_T
    \sim
    \Normal\left(
        \ln 80+\left(0.05-\frac{1}{2}0.25^2\right)0.5,
        0.25^2(0.5)
    \right).
\]
The log-price mean is
\[
    \ln 80+(0.05-0.03125)0.5
    =
    \ln 80+0.009375,
\]
and the log-price standard deviation is
\[
    0.25\sqrt{0.5}=0.1768.
\]
\end{exmp}

\section{Checklist for Students}

After this lecture you should be able to:
\begin{itemize}
    \item distinguish discrete-time, continuous-time, discrete-state, and
    continuous-state stochastic processes;
    \item state and interpret the Markov property;
    \item explain why variance, not standard deviation, adds over independent
    periods;
    \item define a Wiener process and describe its increment distribution;
    \item explain why Brownian paths are continuous but jagged;
    \item write and interpret a generalized Wiener process;
    \item write and interpret an \Ito process;
    \item explain why stock prices are modeled with proportional drift and
    proportional volatility;
    \item simulate one step of geometric Brownian motion;
    \item convert volatility across time units using the square-root-time rule;
    \item construct correlated normal shocks from independent normal shocks;
    \item use the stochastic multiplication rules $(dt)^2=0$, $dt\,dz=0$, and
    $(dz)^2=dt$;
    \item state \Ito's lemma for one stochastic variable;
    \item apply \Ito's lemma to simple functions such as $x^2$, $e^x$, and
    $\ln S$;
    \item derive the process for $\ln S$ when $S$ follows geometric Brownian
    motion;
    \item state the lognormal distribution of $S_T$;
    \item distinguish real-world and risk-neutral stock price dynamics.
\end{itemize}

\section{Practice Questions}

\begin{enumerate}
    \item A process has annual variance rate 4. What are the variance and
    standard deviation of the change over 3 months?
    \item If $z_t$ is a Wiener process and $z_0=10$, what is the distribution
    of $z_2$?
    \item Explain in words why Brownian motion does not have an ordinary
    derivative.
    \item A generalized Wiener process follows $dx=5dt+2dz$. What are the mean,
    variance, and standard deviation of $x_T-x_0$ for $T=3$?
    \item A cash balance follows $dx=12dt+8dz$ and starts at 40. What is the
    distribution after one half year?
    \item A stock follows $dS/S=0.10dt+0.25dz$. What are $\mu$ and $\sigma$?
    Interpret both.
    \item For the stock in the previous question, compute the approximate mean
    and standard deviation of the proportional change over one month.
    \item If annual volatility is $18\%$, what is daily volatility using 252
    trading days?
    \item If daily volatility is $1.8\%$, what is annualized volatility using
    252 trading days?
    \item Simulate one Euler step for $S=100$, $\mu=0.08$, $\sigma=0.20$,
    $\dt=1/12$, and $\eps=-0.75$.
    \item Draw independent $u=0.30$ and $v=1.20$. Construct two standard normal
    shocks with correlation $\rho=0.50$ using the formula in the notes.
    \item Explain why $(dz)^2$ is treated as $dt$ but $dt\,dz$ is treated as
    zero.
    \item State \Ito's lemma for $dx=a(x,t)dt+b(x,t)dz$ and $G=G(x,t)$.
    \item Use \Ito's lemma to find the process for $G=x^3$ when $dx=a\,dt+b\,dz$
    with constant $a$ and $b$.
    \item Use \Ito's lemma to find the process for $G=1/S$ when
    $dS=\mu Sdt+\sigma Sdz$.
    \item Derive $d\ln S$ from $dS/S=\mu dt+\sigma dz$.
    \item If $S_0=100$, $\mu=10\%$, $\sigma=40\%$, and $T=1$, write the
    distribution of $\ln S_T$.
    \item For the previous question, compute the expected stock price and the
    median stock price.
    \item Under risk-neutral dynamics for a non-dividend-paying stock, what
    replaces $\mu$? What happens to $\sigma$?
    \item Explain why the real-world expected return $\mu$ is not the key input
    for Black-Scholes-Merton option pricing.
    \item A derivative value is $G(S,t)$. Explain why the fact that $S$ and $G$
    share the same Brownian shock is essential for delta hedging.
    \item What part of the Lecture 5 material is a bridge to Lecture 6 rather
    than a pricing result by itself?
\end{enumerate}

\end{document}
